\documentclass[12pt]{article}
\usepackage{graphicx} % Required for inserting images
\usepackage{amsmath, amssymb}
\usepackage{geometry}
\usepackage{amsmath, amssymb, amsthm}
\usepackage{geometry}
\geometry{margin=1in}
\usepackage{mathtools}
\usepackage{enumitem}
\usepackage{hyperref}


\title{\textbf{CS687 EndSem Solutions}}
\author{Krishna Kumar Bais\\Roll No.: 241110038}
\date{}

\begin{document}

\maketitle

\section*{Question 1:}

\textbf{Show that if there is a constant \( c \) such that for the sequence \( X \in \Sigma^\infty \), we have}
\[
C(X[0\ldots(n-1)] \mid n) < c \text{ for all } n,
\]
\textbf{then \( X \) is computable.}


\section*{Solution:}
\subsection*{-$>$ Goal:}  
We need to \textbf{prove that \( X \) is computable}, i.e., there exists a total computable function \( f : \mathbb{N} \to \{0,1\} \) such that \( f(n) = X(n) \), the \( n \)th bit of \( X \).

\subsection*{-$>$ Understand the assumption:}

   We are given that:

\[
\forall n \in \mathbb{N}, \quad C(X[0\ldots(n-1)] \mid n) < c,
\]
\newline
This means: for each \( n \), there exists a \textbf{program \( p_n \)} of length \( < c \) such that when run on input \( n \), it outputs \( X[0\ldots(n-1)] \).

So:
\begin{itemize}
  \item \( U(p_n, n) = X[0\ldots(n-1)] \),
  \item where \( U \) is a fixed universal Turing machine.
\end{itemize}

\subsection*{-$>$ There are finitely many such programs:}

Since \( C(\cdot \mid n) < c \), and program size \( < c \), there are at most:

\[
N = 2^{c} - 1
\]
\newline
distinct programs that could serve as \( p_n \) for \textbf{any} \( n \). Let the full list of such programs be:

\[
\mathcal{P} = \{p_1, p_2, \ldots, p_N\}
\]
\newline
Each \( p_i \) is a candidate decoder for \( X[0..n-1] \), when given \( n \).

\subsection*{-$>$ Fix one program that works for all large \( n \):}


Each \( p_i \in \mathcal{P} \) defines a \textbf{partial function} \( f_i(n) = U(p_i, n) \), potentially defined for infinitely many \( n \).
\newline
We know that for each \( n \), \textbf{at least one} \( p_i \) satisfies:

\[
U(p_i, n) = X[0\ldots(n-1)]
\]
\newline
But to conclude \textbf{computability}, we need \textbf{one single program} \( p_i \) that works for \textbf{all sufficiently large} \( n \).

\subsection*{-$>$ Pigeonhole principle → one \( p^* \in \mathcal{P} \) works infinitely often:}

\begin{itemize}
  \item There are only finitely many \( p_i \).
  \item So, by the pigeonhole principle, \textbf{some program} \( p^* \in \mathcal{P} \) must work for \textbf{infinitely many \( n \)}.
\end{itemize}

Let:

\[
S = \{ n \in \mathbb{N} : U(p^*, n) = X[0\ldots(n-1)] \}
\]

Then \( S \subseteq \mathbb{N} \) is infinite.

\subsection*{-$>$ Show that \( S \) must be cofinite (i.e., missing only finitely many n):}

Let us suppose, for contradiction, that \( S \) is \textbf{not cofinite}; i.e., \( \mathbb{N} \setminus S \)  is infinite.
\newline
\newline
Then infinitely many \( n \) use programs \textbf{different} from \( p^* \). But again, the number of programs is finite, so some other \( p_i \) would also be used infinitely often.
\newline
\newline
This process leads to an infinite sequence of switches between candidate programs — which would imply \textbf{infinitely many distinct infinite sequences} are being reconstructed by programs of size \( < c \).
\newline
\newline
But there are \textbf{only finitely many} such programs, and each program \( p \) defines at most \textbf{one infinite sequence} (since \( U(p, n) \) must agree with \( X[0\ldots(n-1)] \) for all large \( n \)).
\newline
Thus, we reach a contradiction unless \textbf{one \( p^* \)} outputs \( X[0..n-1] \) for \textbf{all sufficiently large \( n \)}.
\newline
\newline
So there exists a \textbf{fixed program \( p^* \)} such that:

\[
\exists N_0\ \forall n \geq N_0,\quad U(p^*, n) = X[0\ldots(n-1)]
\]

\subsection*{-$>$ Use \( p^* \) to define a total computable function \( f(n) = X(n) \):}

Define the function \( f: \mathbb{N} \to \{0,1\} \) as:

\begin{itemize}
  \item For \( n < N_0 \), hardcode the value of \( X(n) \).
  \item For \( n \geq N_0 \):  
  \begin{enumerate}
    \item Compute \( X[0\ldots(n-1)] = U(p^*, n) \)
    \item Return \( X(n) = \text{nth bit of } U(p^*, n) \)
  \end{enumerate}
\end{itemize}

Since:
\begin{itemize}
  \item \( p^* \) is fixed and finite,
  \item \( U \) is universal (and total on this domain),
  \item and the hardcoded part is finite,
\end{itemize}

\(\Rightarrow\ f(n) \) is total and computable.

\subsection*{Conclusion:}
\[
X(n) = f(n),\quad \text{where } f \text{ is total computable}
\]
Hence, the infinite sequence \( X \) is computable.
\newline
\newline
\textit{If there is a constant \( c \) such that for all \( n \), \( C(X[0\ldots(n-1)] \mid n) < c \),\\
then \( X \) is computable.}

\newpage
\section*{Question 2:}

\noindent
\textbf{Show that if for any infinite binary sequence \( X \):}  
\[
\sum_{n \in \mathbb{N}} 2^{n - K(X[0\ldots(n-1)])} < \infty \tag{1}
\]
\textbf{then:}  
\[
\lim_{n \to \infty} K(X[0 \ldots (n-1)]) - n = \infty \tag{2}
\]

\section*{Solution:}
\subsection*{-$>$ Goal:}

We need to prove that if the sum in (1) is finite, then the quantity \( K(X[0\ldots(n-1)]) - n \) grows unbounded — i.e., the prefix complexities grow \textbf{strictly faster than linear}, asymptotically.

\subsection*{-$>$ Intuition and Structure:}

We are asked to show that the prefix complexity \( K(X[0\ldots(n-1)]) \) must eventually grow faster than \( n \), under the assumption that the sum  
\[
\sum_{n \in \mathbb{N}} 2^{n - K(X[0\ldots(n-1)])}
\]
is finite.
\newline
\newline
Let us define:
\begin{itemize}
  \item \( s_n := K(X[0\ldots(n-1)]) \)
\end{itemize}

So (1) becomes:
\[
\sum_{n \in \mathbb{N}} 2^{n - s_n} < \infty
\]

We want to conclude that:
\[
\lim_{n \to \infty} (s_n - n) = \infty
\quad \text{(i.e., } s_n \gg n \text{)}
\]

If the difference \( s_n - n \) were bounded, say \( s_n \leq n + c \), then the sum becomes:
\[
\sum_n 2^{n - s_n} \geq \sum_n 2^{n - (n + c)} = \sum_n 2^{-c} = \infty
\]
— contradiction.


\subsection*{-$>$ Assumption (for contradiction):}

Suppose the limit in (2) does \textbf{not} go to \( \infty \). Then there exists a constant \( c \in \mathbb{N} \) and an \textbf{infinite} subsequence \( \{n_k\} \) such that:

\[
K(X[0\ldots(n_k - 1)]) \leq n_k + c
\quad \text{for infinitely many } k
\]

Then:

\[
2^{n_k - K(X[0\ldots(n_k - 1)])} \geq 2^{n_k - (n_k + c)} = 2^{-c}
\]

Therefore, for infinitely many \( n_k \):

\[
2^{n_k - K(X[0\ldots(n_k - 1)])} \geq 2^{-c}
\]

So the sum:

\[
\sum_n 2^{n - K(X[0\ldots(n-1)])}
\]

has infinitely many terms each \( \geq 2^{-c} \), hence:

\[
\sum_n 2^{n - K(X[0\ldots(n-1)])} = \infty
\]

This \textbf{contradicts the given assumption} (1) that the sum is finite.



\subsection*{Conclusion:}

Therefore, our assumption must be false, and the only possibility is:

\[
\lim_{n \to \infty} K(X[0\ldots(n-1)]) - n = \infty \tag{\textbf{proved}}
\]
\newline
\newline
If \(\sum_{n \in \mathbb{N}} 2^{n - K(X[0\ldots(n-1)])} < \infty\) for an infinite binary sequence \(X\), then:  
\[
\boxed{\lim_{n \to \infty} K(X[0\ldots(n-1)]) - n = \infty}
\]


\newpage
\section*{Question 3:}

Suppose \( d : \Sigma^* \to [0, \infty) \) is a \textbf{lower semicomputable martingale}. Let \( X \) be a sequence on which \(d\) \textbf{succeeds}.  
Then show that \( X \) has \textbf{infinitely many prefixes with low self-delimiting Kolmogorov complexity}.

\section*{Solution:}
\subsection*{-$>$ Goal:}

We need to prove:

\[
\text{If } \limsup_{n \to \infty} d(X[0 \ldots n-1]) = \infty, \text{ then}
\]
\[
\exists \text{ infinitely many } \sigma \prec X \text{ such that } K(\sigma) \leq |\sigma| - \log d(\sigma) + O(1)
\]



\subsection*{-$>$ Intuition:}

A \textbf{lower semicomputable martingale} \( d \) is a computable betting strategy. If it \textbf{succeeds} on an infinite binary sequence \( X \), that means \( d(X[0..n]) \to \infty \) along a subsequence — i.e., it wins unbounded money by betting on \( X \).
\newline
\newline
Such sequences must then \textbf{fail randomness tests} (they are not Martin-Löf random). One way to witness failure of randomness is that \textbf{some prefixes of the sequence have unexpectedly low prefix-free Kolmogorov complexity} — they are "too simple to be random".
\newline
\newline
We’ll formalize this using \textbf{Kraft's inequality}, \textbf{code construction}, and known theorems (e.g., from class notes and Li \& Vitányi) connecting martingale success to complexity bounds..



\subsubsection*{-$>$ Definitions:}

Let:
\begin{itemize}
  \item \( \Sigma^* \) = set of finite binary strings.
  \item \( d : \Sigma^* \to \mathbb{R}_{\geq 0} \) is a \textbf{lower semicomputable martingale}:
    \begin{itemize}
      \item \( d(\sigma) = \frac{1}{2}(d(\sigma 0) + d(\sigma 1)) \)
      \item Computably approximable from below.
    \end{itemize}
  \item \( d \) \textbf{succeeds} on \( X \in \Sigma^\infty \):
  \[
  \limsup_{n \to \infty} d(X[0..n-1]) = \infty
  \]
\end{itemize}

Let \( \sigma_n = X[0..n-1] \). We aim to show that:

\[
\text{For infinitely many } n, \quad K(\sigma_n) \leq |\sigma_n| - \log d(\sigma_n) + O(1)
\]

\subsubsection*{-$>$ Use the Kraft Inequality via Ville’s Lemma:}

We define a set of strings where the martingale exceeds some fixed capital:

Let \( c \in \mathbb{N} \). Define:

\[
S_c = \{ \sigma \in \Sigma^* : d(\sigma) \geq 2^c \}
\]

Ville’s inequality gives:

\[
\sum_{\sigma \in S_c} 2^{-|\sigma|} \leq 2^{-c}
\]

This allows us to define a \textbf{prefix-free code} for strings in \( S_c \), with lengths:

\[
\ell(\sigma) \approx |\sigma| - \log d(\sigma) + O(1)
\]

To explain this: Since \( \sum_{\sigma} 2^{-\ell(\sigma)} \leq 1 \), such codes obey \textbf{Kraft's inequality} and can be interpreted as the lengths of programs (i.e., descriptions) in a prefix code.

\subsubsection*{-$>$ Prefix Complexity from Code Lengths:}

Using the fact that the prefix-free Kolmogorov complexity \( K(\sigma) \) is upper-bounded by the length of a prefix code for \( \sigma \):

\[
K(\sigma) \leq \ell(\sigma) + O(1) \leq |\sigma| - \log d(\sigma) + O(1)
\]

This bound is \textbf{tight and valid} for all \( \sigma \in S_c \).

\subsubsection*{-$>$ Apply to \( X \):}

Since \( d \) succeeds on \( X \), for \textbf{infinitely many} \( n \), we have:

\[
d(X[0..n-1]) \geq 2^c \quad \text{for arbitrarily large } c
\]

So:
\begin{itemize}
  \item For infinitely many \( n \), \( \sigma_n \in S_c \) for increasing \( c \)
  \item For such \( \sigma_n \), we have:
  \[
  K(\sigma_n) \leq |\sigma_n| - \log d(\sigma_n) + O(1)
  \]
\end{itemize}

This proves that \( X \) has \textbf{infinitely many prefixes} with \textbf{lower-than-expected Kolmogorov complexity}, compared to the length of the prefix.

\subsection*{Conclusion:}

If \( d \) is a lower semicomputable martingale and \( d \) succeeds on an infinite binary sequence \( X \), then:

\[
\text{There are infinitely many } \sigma \prec X \text{ such that:} \quad
K(\sigma) \leq |\sigma| - \log d(\sigma) + O(1)
\]

i.e., \( X \) has infinitely many prefixes with \textbf{low prefix-free Kolmogorov complexity}.  

\newpage
\section*{Problem 4:}
 \textbf{Show that the following inequality holds for all strings $x, y$ and $z$ of length $n$ .}
\[
2C(x, y, z) \leq C(x, y) + C(x, z) + C(y, z) + O(\log n)
\]

\section*{Solution:}

\subsection*{-$>$ Preliminaries and Definitions:}

Let $C(u)$ denote the plain Kolmogorov complexity of string $u$. \\ 
Let $C(u, v)$ denote the complexity of the pair $(u, v)$, i.e., the length of the shortest program that outputs both $u$ and $v$.
\newline
\newline
Standard facts (used without proof here):

\begin{enumerate}
    \item \textbf{Chain rule}:
    \[
    C(u, v) = C(u) + C(v|u) + O(\log n)
    \]

    \item \textbf{Triple chain rule}:
    \[
    C(u, v, w) = C(u, v) + C(w | u, v) + O(\log n)
    \]

    \item \textbf{Symmetry of information}:
    \[
    C(v|u) + C(u) = C(u, v) + O(\log n)
    \]

    \item For any strings $u, v, w$:
    \[
    C(v|u, w) + C(w|u, v) \leq C(v, w|u) + O(\log n)
    \]
\end{enumerate}

\subsection*{-$>$ Step-by-step Proof:}

\subsubsection*{Step 1: Write $C(x, y, z)$ in two ways}

Use the triple chain rule:

\[
C(x, y, z) = C(x, y) + C(z | x, y) + O(\log n) \tag{1}
\]

\[
C(x, y, z) = C(x, z) + C(y | x, z) + O(\log n) \tag{2}
\]

\subsubsection*{Step 2: Add the two expressions}

Add equations (1) and (2):

\[
2C(x, y, z) = C(x, y) + C(x, z) + C(z | x, y) + C(y | x, z) + O(\log n) \tag{3}
\]

\subsubsection*{Step 3: Bound the conditional complexities}

Use the inequality:

\[
C(z | x, y) + C(y | x, z) \leq C(y, z | x) + O(\log n) \tag{4}
\]

Then add this to (3):

\[
2C(x, y, z) \leq C(x, y) + C(x, z) + C(y, z | x) + O(\log n) \tag{5}
\]

Now use the chain rule again:

\[
C(y, z | x) \leq C(y, z) - C(x) + O(\log n) \tag{6}
\]

But this doesn't simplify directly. Instead, we use the looser inequality:

\[
C(y, z | x) \leq C(y, z) + O(\log n)
\]

Substitute into (5):

\[
2C(x, y, z) \leq C(x, y) + C(x, z) + C(y, z) + O(\log n) \tag{7}
\]

\subsection*{Conclusion:}

We have derived:

\[
2C(x, y, z) \leq C(x, y) + C(x, z) + C(y, z) + O(\log n)
\]
\\
which holds for all strings $x, y, z$ of length at most $n$, and the hidden constant in $O(\log n)$ is independent of the strings.



\end{document}
